\tzpanchor{1005}
\tzpentryheader{1005}{BESSEL-WEIGHTED TIDAL TORQUE}

\begin{tzpsnapshot}
\tzpsnaprow{TZPID}{\tzpcode{ID1005}}
\tzpsnaprow{UUID}{\tzpcode{fda208f9-ea54-577c-8c2d-de8c935e7fbf}}
\tzpsnaprow{Type}{Transfer Statement}
\tzpsnaprow{Scale}{transfer}
\tzpsnaprow{LOC Classification}{Working routing: QA641 manifolds and topology; QC174.17 mathematical physics}
\tzpsnaprow{arXiv Classification}{Primary: math-ph; Cross-list: gr-qc, math.DG}
\end{tzpsnapshot}

\tzpsection{Canonical Statement}
tau\_\{tidal\}(omega, Omega) = frac\{1\}\{I\} sum\_\{n=1\}\textasciicircum{}\{N\} J\_n(ne)\textasciicircum{}2 sin(2(omega - nOmega))

\[
dot\varphi_i = \omega_i + \frac{1}{N}\sum_j K_{ij}\sin(\varphi_j-\varphi_i) + \eta_i(t).
\]
\[
lambda_i(t) = \lambda_0 \sum_n J_n(\kappa)^2 S_{i,n}(t) [1+\beta R_i(t)].
\]

\noindent
\begin{minipage}[t]{0.48\linewidth}
\tzpsection{Wolfram Form}
\begingroup
\small
\[
\begin{aligned}
\mathfrak{W}_{\mathrm{TZP}}\!\left(\mathcal{K}_{1005}\right)
&\longmapsto
\left\langle \Omega^{\mathbb{T}},\; \Psi_{\mathrm{T}},\; \rho_{\mathrm{vac}} \right\rangle
\end{aligned}
\]
\[
\mathcal{K}_{1005}
:=
\operatorname{HoldForm}\!\left(\mathcal{T}_{1005}\right)
\]
\[
\begin{aligned}
\operatorname{Admissible}_{\mathrm{TZP}}\!\left(\mathcal{K}_{1005}\right)
&\Longleftrightarrow
\mathfrak{W}_{\mathrm{TZP}}\!\left(\mathcal{K}_{1005}\right)\not\equiv\varnothing
\end{aligned}
\]
\textbf{Wolfram register:} canonical symbol lift\\
\textbf{Title lift:} BESSEL-WEIGHTED TIDAL TORQUE\\
\textbf{Symbol key:} \(\Omega^{\mathbb{T}}\): Trawinistic boundary curvature\\ \(\Psi_{\mathrm{T}}\): Trawinistic wavefunction\\ \(\rho_{\mathrm{vac}}\): vacuum energy density
\par
\endgroup
\end{minipage}\hfill
\begin{minipage}[t]{0.48\linewidth}
\tzpsection{TRAWIN Normalization}
\[
\tzpnormalizewrap{dot\varphi_i = \omega_i + \frac{1}{N}\sum_j K_{ij}\sin(\varphi_j-\varphi_i) + \eta_i(t).,\,lambda_i(t) = \lambda_0 \sum_n J_n(\kappa)^2 S_{i,n}(t) [1+\beta R_i(t)].,\,tau_{sync} \approx \frac{\omega I Q}{3 k_2 G M_P^2 R^5} a^6}
\]

\small
\textbf{T}: Mode Quantization is localized within the category theory and proof structure arc of the directory.\\
\textbf{R}: Dispersion Law preserves the operative relation that bessel-weighted tidal torque requires.\\
\textbf{A}: Boundary Condition fixes the admissible closure condition for the entry.\\
\textbf{W}: BESSEL-WEIGHTED TIDAL TORQUE is rendered as a reusable operator-level statement.\\
\textbf{I}: The informational content of bessel-weighted tidal torque is retained rather than flattened.\\
\textbf{N}: BESSEL-WEIGHTED TIDAL TORQUE closes as a distinct normalized TRAWIN object.
\normalsize
\end{minipage}

\tzpsection{Derivable Consequences}
The entry is now carried by explicit resonance mathematics, which better matches the wave-and-cymatic story arc of the third hundred.
\clearpage

\tzpentryheader{1005}{Oxford Dictionary of the Queens, NY English}

\begingroup\small
\tzpsection{Definition}
BESSEL-WEIGHTED TIDAL TORQUE is a registry definition in Category Theory and Proof Structure with secondary pressure from Signal.

% BEGIN TZP CONTEXTUAL MEANING
\tzpsection{Contextual Meaning}
\begingroup\footnotesize\setlength{\emergencystretch}{3em}\sloppy
\textbf{Formal statement:} tau sub tidal (omega, Omega) = frac 1 I sum sub n=1 to the power N J sub n(ne) to the power 2 sin(2(omega - nOmega)). \textbf{Contextual meaning:} The entry reads BESSEL-WEIGHTED TIDAL TORQUE as a controlled technical headword whose equation layer, proof route, and registry role are interpreted together. \textbf{Derivable consequences:} The entry is now carried by explicit resonance mathematics, which better matches the wave-and-cymatic story arc of the third hundred. \textbf{Lemma support:} Mode Quantization; Dispersion Law; Boundary Condition.\par
\endgroup
% END TZP CONTEXTUAL MEANING

\tzpsection{Technical Interpretation}
Technically, BESSEL-WEIGHTED TIDAL TORQUE is read through the local equation chain and the TRAWIN normalization recorded on the entry page.

\tzpsection{Equation Grounding}
Equation anchors: dotvarphi sub i = omega sub i + ratio 1 N sum sub j K sub ij sin(varphi sub j-varphi sub i) + eta sub i(t); lambda sub i(t) = lambda sub 0 sum sub n J sub n(kappa) to the power 2 S sub i,n (t) [1+beta R sub i(t)]; and tau sub sync approximately ratio omega I Q 3 k sub 2 G M sub P to the power 2 R to the power 5 a to the power 6. Proof route: show that bounded carrier geometry forces the stated mode quantization and resonance law. Formalization route: prove that the stated boundary condition yields the stated discrete spectrum.

\tzpsection{Interpretive Role}
The entry is now carried by explicit resonance mathematics, which better matches the wave-and-cymatic story arc of the third hundred.
\endgroup

\tzpsection{TZP Thesaurus}
\begin{enumerate}[leftmargin=1.6em,itemsep=6pt,topsep=4pt]
\item \textbf{Cylindrical Mode Weighted Tidal Torque}: nearby spectral language for modes, resonant structure, and admissible oscillation.
\item \textbf{Cylindrical Mode Spectral Analogue}: nearby spectral language for modes, resonant structure, and admissible oscillation.
\item \textbf{Cylindrical Mode Terminology}: nearby spectral language for modes, resonant structure, and admissible oscillation.
\end{enumerate}

\tzpsection{Encyclopedia Mathematica}
\begin{samepage}
\noindent\begin{minipage}[t]{0.45\linewidth}
\vspace{0pt}
\raggedright
\scriptsize\textbf{Image reading.} The rendered manifold at right is the per-ID light plate for BESSEL-WEIGHTED TIDAL TORQUE. It should be read as the visual companion to the entry's equation layer, not as a repeated template diagram.\par
\end{minipage}\hfill
\begin{minipage}[t]{0.53\linewidth}
\vspace*{-0.10in}
\raggedright
\tzpencyclopediaimage{1005}
\end{minipage}
\end{samepage}

\clearpage

\tzpentryheader{1005}{Direct Equation Progression and Proof Trajectory}

\tzpsection{Direct Equation Progression}
\begin{enumerate}[leftmargin=1.6em,itemsep=9pt,topsep=6pt]
\item \textbf{Mode Quantization}
\[
dot\varphi_i = \omega_i + \frac{1}{N}\sum_j K_{ij}\sin(\varphi_j-\varphi_i) + \eta_i(t).
\]
This fixes the quantized carrier mode indexed by the bounded geometry.

\item \textbf{Dispersion Law}
\[
lambda_i(t) = \lambda_0 \sum_n J_n(\kappa)^2 S_{i,n}(t) [1+\beta R_i(t)].
\]
This turns the mode index into an actual propagation or resonance law.

\item \textbf{Boundary Condition}
\[
tau_{sync} \approx \frac{\omega I Q}{3 k_2 G M_P^2 R^5} a^6
\]
This closes the progression by enforcing the boundary or nodal condition that makes the pattern admissible.
\end{enumerate}

\tzpsection{Proof}
\textbf{Proof route.} show that bounded carrier geometry forces the stated mode quantization and resonance law.

\textbf{Formal method.} prove that the stated boundary condition yields the stated discrete spectrum.

\medskip
\tzpproofartifacts{Isabelle audit: corpus classification recorded in appendix.}{Lean audit: compiled corpus certificate.}{Rocq audit: compiled corpus certificate.}

\tzpsection{Dependencies and Test Handle}
\begin{tzpsnapshot}
\tzpsnaprow{Dependencies}{\tzpidref{1004}, \tzpidref{1000}}
\tzpsnaprow{Node}{\tzpcode{registry_id_1005}}
\tzpsnaprow{Support Titles}{Mode Quantization; Dispersion Law; Boundary Condition}
\tzpsnaprow{Test Handle}{If the resonance pattern cannot be organized into bounded modes with the stated boundary condition, the entry fails.}
\end{tzpsnapshot}

\tzpsection{Canonical Equation Links}
\tzpequationlinks
  {k_{\mathrm{n}} = \frac{n\pi}{L}}
  {\omega_{\mathrm{n}} = v\,k_{\mathrm{n}}}
  {\Phi_{\mathrm{n}}|_{\partial\Omega}=0}
